\documentclass[12pt]{ximera}
\title{Exam 1 Review}
\begin{document}
\begin{abstract}
Review for Exam 1, covering Chapter 1: full rankings by all four voting methods,
Condorcet candidates and Condorcet losers, an Independence-of-Irrelevant-Alternatives
violation under the Borda count, counting pairwise comparisons, and a modified Borda
count and ``most last-place votes'' elimination.
\end{abstract}
\maketitle

\section*{Exam 1 Review}

Topics: voting systems, Plurality, Borda count, Plurality-with-Elimination, Pairwise
Comparisons, and the fairness criteria with Arrow's Impossibility Theorem. Once you are
comfortable with the material, try these with a stopwatch: you should be able to finish
them in under $60$ minutes.

\begin{problem}
Use the following preference schedule to find the full ranking of candidates A, B, C,
D, E by each method.
\begin{center}
\begin{tabular}{c|ccccc}
Number of votes & 20 & 15 & 7 & 3 & 2\\\hline
1st & A & C & E & B & D\\
2nd & B & E & D & E & C\\
3rd & C & B & B & A & B\\
4th & D & D & C & C & A\\
5th & E & A & A & D & E
\end{tabular}
\end{center}

\begin{prompt}
\textbf{(a)} Plurality: 1st $\answer{A}$, 2nd $\answer{C}$, 3rd $\answer{E}$,
4th $\answer{B}$, 5th $\answer{D}$

\begin{feedback}
First-place votes: A $20$, C $15$, E $7$, B $3$, D $2$.
\end{feedback}
\end{prompt}

\begin{prompt}
\textbf{(b)} Borda count: 1st $\answer{B}$, 2nd $\answer{C}$, 3rd $\answer{A}$,
4th $\answer{E}$, 5th $\answer{D}$

\begin{feedback}
With $5,4,3,2,1$ points: B $167$, C $163$, A $135$, E $129$, D $111$. (Check: $47$
ballots at $15$ points each is $705$, and $167+163+135+129+111=705$.)
\end{feedback}
\end{prompt}

\begin{prompt}
\textbf{(c)} Plurality-with-Elimination: 1st $\answer{C}$, 2nd $\answer{A}$,
3rd $\answer{E}$, 4th $\answer{B}$, 5th $\answer{D}$

\begin{feedback}
\begin{itemize}
\item Round 1: A $20$, C $15$, E $7$, B $3$, D $2$. Eliminate D; its $2$ votes go to C.
\item Round 2: A $20$, C $17$, E $7$, B $3$. Eliminate B; its $3$ votes go to E.
\item Round 3: A $20$, C $17$, E $10$. Eliminate E; its $7$ original voters rank C
      above A, and the $3$ former B voters rank A above C.
\item Round 4: C $24$, A $23$. C wins.
\end{itemize}
\end{feedback}
\end{prompt}

\begin{prompt}
\textbf{(d)} Pairwise Comparisons: 1st $\answer{B}$, 2nd $\answer{C}$, 3rd $\answer{E}$,
4th $\answer{D}$, 5th $\answer{A}$

\begin{feedback}
\[
\begin{array}{lll@{\qquad}lll}
\text{A vs B} & 20\text{--}27 & \text{B} & \text{B vs D} & 38\text{--}9 & \text{B}\\
\text{A vs C} & 23\text{--}24 & \text{C} & \text{B vs E} & 25\text{--}22 & \text{B}\\
\text{A vs D} & 23\text{--}24 & \text{D} & \text{C vs D} & 38\text{--}9 & \text{C}\\
\text{A vs E} & 22\text{--}25 & \text{E} & \text{C vs E} & 37\text{--}10 & \text{C}\\
\text{B vs C} & 30\text{--}17 & \text{B} & \text{D vs E} & 22\text{--}25 & \text{E}
\end{array}
\]
Points: B $4$, C $3$, E $2$, D $1$, A $0$.
\end{feedback}
\end{prompt}
\end{problem}

\begin{problem}
Is there a Condorcet candidate for the election in Problem 1?

\begin{multipleChoice}
\choice{Yes, A}
\choice[correct]{Yes, B}
\choice{Yes, C}
\choice{Yes, E}
\choice{No}
\end{multipleChoice}

\begin{feedback}
Yes: B wins all four of its head-to-head comparisons. This is why B also wins the
Pairwise Comparisons method.
\end{feedback}
\end{problem}

\begin{problem}
A candidate who loses in pairwise comparison against every other candidate is called a
\emph{Condorcet loser}. Is there a Condorcet loser in the election in Problem 1?

\begin{multipleChoice}
\choice[correct]{Yes, A}
\choice{Yes, D}
\choice{Yes, E}
\choice{No}
\end{multipleChoice}

\begin{feedback}
Yes: A loses every comparison, some by as little as $23$--$24$. A is the Plurality
winner and at the same time the Condorcet loser.
\end{feedback}
\end{problem}

\begin{problem}
Consider the following preference schedule.
\begin{center}
\begin{tabular}{c|ccc}
Number of votes & 10 & 8 & 3\\\hline
1st & B & C & A\\
2nd & A & B & C\\
3rd & C & A & B
\end{tabular}
\end{center}

\begin{prompt}
\textbf{(a)} Find the Borda points and the winner.

A: $\answer{37}$ \quad B: $\answer{49}$ \quad C: $\answer{40}$ \qquad Winner:
$\answer{B}$

\begin{feedback}
\[
\begin{aligned}
\text{A:}&\ 10\cdot2+8\cdot1+3\cdot3=37\\
\text{B:}&\ 10\cdot3+8\cdot2+3\cdot1=49\\
\text{C:}&\ 10\cdot1+8\cdot3+3\cdot2=40
\end{aligned}
\]
\end{feedback}
\end{prompt}

\begin{prompt}
\textbf{(b)} If A withdraws, who wins by the Borda count? $\answer{C}$

\begin{feedback}
With only B and C, the points are B: $10\cdot2+8\cdot1+3\cdot1=31$ and C:
$10\cdot1+8\cdot2+3\cdot2=32$. C wins. A non-winner's withdrawal changed the winner, so
the Independence-of-Irrelevant-Alternatives criterion is violated.
\end{feedback}
\end{prompt}

\begin{prompt}
\textbf{(c)} Which other fairness criteria are violated in this example?

\begin{selectAll}
\choice{Majority criterion}
\choice{Condorcet criterion}
\choice{Monotonicity criterion}
\choice[correct]{None of these}
\end{selectAll}

\begin{feedback}
None. No candidate has a majority ($11$ of $21$) of first-place votes. There is no
Condorcet candidate: B beats A $18$--$3$, A beats C $13$--$8$, and C beats B $11$--$10$,
which is a cycle. The Borda count never violates monotonicity.
\end{feedback}
\end{prompt}
\end{problem}

\begin{problem}
In an election with $200$ candidates using the Pairwise Comparisons system, how many
pairwise comparisons must be evaluated? $\answer{19900}$

\begin{feedback}
With $N$ candidates there are $\frac{N(N-1)}{2}$ comparisons:
$\frac{200\cdot199}{2}=19{,}900$.
\end{feedback}
\end{problem}

\begin{problem}
Below is the preference schedule for an election for Best Ice Cream Flavor. The candidates are
Almonds \& Joy, Butterscotch, Candycorn, and Drizzle Deluxe (A, B, C, D).
\begin{center}
\begin{tabular}{c|ccccc}
Number of votes & 10 & 7 & 6 & 5 & 2\\\hline
1st & A & B & C & C & A\\
2nd & B & D & B & D & C\\
3rd & C & A & D & B & D\\
4th & D & C & A & A & B
\end{tabular}
\end{center}

\begin{prompt}
\textbf{(a)} Find the points and the winner using a modified Borda count with $5$
points for first place, $3$ for second, $1$ for third, and $0$ for fourth.

A: $\answer{67}$ \quad B: $\answer{88}$ \quad C: $\answer{71}$ \quad D: $\answer{44}$
\qquad Winner: $\answer{B}$

\begin{feedback}
\[
\begin{aligned}
\text{A:}&\ 10\cdot5+7\cdot1+2\cdot5=67\\
\text{B:}&\ 10\cdot3+7\cdot5+6\cdot3+5\cdot1=88\\
\text{C:}&\ 10\cdot1+6\cdot5+5\cdot5+2\cdot3=71\\
\text{D:}&\ 7\cdot3+6\cdot1+5\cdot3+2\cdot1=44
\end{aligned}
\]
(Check: $30$ ballots at $9$ points each is $270=67+88+71+44$.) Butterscotch wins.
\end{feedback}
\end{prompt}

\begin{prompt}
\textbf{(b)} Now use a modified Plurality-with-Elimination, where each round eliminates
the candidate with the \emph{most last-place} votes. A ballot whose last-place candidate
is eliminated counts toward its next-worst candidate. The last candidate standing wins,
the next-to-last is second, and so on.

Round 1 last-place votes: A $\answer{11}$, B $\answer{2}$, C $\answer{7}$,
D $\answer{10}$

\smallskip
Ranking: 1st $\answer{B}$, 2nd $\answer{C}$, 3rd $\answer{D}$, 4th $\answer{A}$

\begin{feedback}
\begin{itemize}
\item Round 1: A $11$, D $10$, C $7$, B $2$. Eliminate A.
\item Round 2: the $6+5$ ballots with A last now count for D and B, so D $16$,
      B $7$, C $7$. Eliminate D.
\item Round 3: B $13$, C $17$. Eliminate C, leaving B.
\end{itemize}
The elimination order A, D, C reverses to the ranking B, C, D, A.
\end{feedback}
\end{prompt}
\end{problem}

\begin{problem}
An election with four candidates and $25$ voters uses the modified Borda count with
$5,3,1,0$ points from the preceding problem. Candidates A, B, C have $62$, $33$, and
$48$ points. How many points did D get? $\answer{82}$

\begin{hint}
Think about the total number of Borda points.
\end{hint}

\begin{feedback}
Each ballot awards $5+3+1+0=9$ points, so the total is $9\cdot25=225$. D has
$225-(62+33+48)=82$.
\end{feedback}
\end{problem}

\begin{problem}
Suppose the Ice Cream Flavor election had produced this schedule instead:
\begin{center}
\begin{tabular}{c|ccccc}
Number of votes & 10 & 7 & 6 & 5 & 2\\\hline
1st & A & B & C & C & A\\
2nd & C & D & B & D & C\\
3rd & B & A & D & B & D\\
4th & D & C & A & A & B
\end{tabular}
\end{center}
Then D withdraws before the winners are announced.

\begin{prompt}
\textbf{(a)} Using the ``most last-place votes'' elimination, who wins with D in the
race? $\answer{C}$

\begin{feedback}
Round 1: A $11$, D $10$, C $7$, B $2$, so A is eliminated. Round 2: D $16$, B $7$,
C $7$, so D is eliminated. Round 3: B $23$, C $7$, so B is eliminated. C wins.
\end{feedback}
\end{prompt}

\begin{prompt}
\textbf{(b)} Who wins after D withdraws? $\answer{A}$

\begin{feedback}
Without D, B moves to last place on the $10$ ballots in the first column. Round 1 then
gives B $10+2=12$, A $11$, C $7$, so B is eliminated. Round 2 gives C $10+7+2=19$ and
A $11$, so C is eliminated and A wins.
\end{feedback}
\end{prompt}

\begin{prompt}
\textbf{(c)} Which fairness criterion does this show the method violates?

\begin{multipleChoice}
\choice{Majority criterion}
\choice{Condorcet criterion}
\choice{Monotonicity criterion}
\choice[correct]{Independence-of-Irrelevant-Alternatives criterion}
\end{multipleChoice}

\begin{feedback}
The withdrawal of D, a non-winner, changed the winner from C to A, so the
Independence-of-Irrelevant-Alternatives criterion is violated.
\end{feedback}
\end{prompt}
\end{problem}

\end{document}
